\chapter*{Resumen}

Este Trabajo de Fin de Grado surge como respuesta a la necesidad que había en el departamento de Ingeniería Aeroespacial y Mecánica de Fluidos de una herramienta que facilitase el análisis de información obtenida en los experimentos de inyección de fluidos. En concreto, se buscaba una solución que permitiera detectar y analizar microgotas en vídeos de alta velocidad, para obtener diversas mediciones como el tamaño, la velocidad o la trayectoria de las gotas. Esta tarea se realizaba de forma manual, lo que resultaba tedioso y propenso a errores, especialmente cuando se trataba de grandes volúmenes de datos.

Para dar respuesta a esta necesidad, se ha desarrollado una aplicación de escritorio utilizando diversos algoritmos de \textbf{visión por computador}, usando las implementaciones disponibles en la biblioteca \textbf{OpenCV\cite{opencv_library}}. Para realizar esta detección se preprocesan los frames para ajustar el contraste y eliminar el ruido, y posteriormente se aplican técnicas de segmentación para identificar las regiones que corresponden a las gotas y operaciones morfológicas para mejorar la calidad de la detección. Una vez detectadas, se extraen características como el tamaño, la velocidad y la trayectoria y se guardan en una estructura de datos para su posterior análisis.

Debido a que los investigadores no tenían conocimientos de programación, se ha diseñado una interfaz gráfica de usuario intuitiva y fácil de usar, que permite a los usuarios cargar vídeos, configurar los parámetros de detección y visualizar los resultados de forma clara y accesible usando \textbf{Tkinter\cite{python_tkinter}}. Además, se han implementado funcionalidades para exportar los datos obtenidos en formato CSV.